\chapter{213}
\label{ch:213}

% Joint research by Mingu Jeong and Claude (Anthropic)

This chapter presents a structure that precedes
the mathematical framework developed in previous chapters.
The results here do not depend on any prior chapter.

\section{The Structure}

Let $E = \{e_1, e_2, e_3\}$ be a set with three elements.
Define two binary operations on formal expressions over $E$:
\begin{itemize}
\item $a + b$: concatenation with $e_1$ between.
\item $a \times b$: concatenation without separator.
\end{itemize}
Define $0 := e_1 + e_1$.

An equivalence $\approx$ is generated by twelve rules.
Nine compensate for the ordering imposed by linear notation
(commutativity, associativity, congruence, reflexivity,
symmetry, transitivity).
Three are intrinsic:
\begin{align}
e_1 + a &\approx a \label{eq:a8} \\
e_1 \times a &\approx e_1 \label{eq:a9} \\
a \times (b + c) &\approx
  (a \times b) + (a \times c) \label{eq:a12}
\end{align}

\section{Swap Symmetry}

The involution $\sigma(e_1) = e_1$,
$\sigma(e_2) = e_3$, $\sigma(e_3) = e_2$
preserves $\approx$: if $a \approx b$ then
$\sigma(a) \approx \sigma(b)$.

This means $e_2$ and $e_3$ are indistinguishable
in the axiom system.
All three intrinsic axioms refer only to $e_1$.
The structure is $\{$boundary, content, content$\}$.

\section{Independence}

The three intrinsic axioms are pairwise independent.
Three counterexample models:
\begin{center}
\begin{tabular}{lllll}
\hline
Model & $S$ & $+$ & $\times$ & Fails \\
\hline
1 & $\{0,1\}$ & max & const $1$ & \eqref{eq:a8} \\
2 & $\{0,1\}$ & max & max & \eqref{eq:a9} \\
3 & $\{z,a,b\}$ & max & custom & \eqref{eq:a12} \\
\hline
\end{tabular}
\end{center}

\section{Fixed Point}

$\binom{n}{2} = n$ has exactly two solutions:
$n = 0$ (trivial) and $n = 3$.

Under iteration of $n \mapsto \binom{n}{2}$:
$n < 3$ converges to $0$;
$n = 3$ is stationary;
$n > 3$ diverges.

\section{Objects and Morphisms}

Let $\mathrm{Obj} = \{a, b, c\}$ and let
$\mathrm{Mor}$ be the set of unordered pairs.
$|\mathrm{Mor}| = \binom{3}{2} = 3 = |\mathrm{Obj}|$.
An explicit bijection exists between the two sets.

Objects and morphisms are the same set.
The number $2$ (arity) and $3$ (cardinality)
are two readings of one structure.

\section{Decidability}

$\approx$ is decidable.
Each expression has a normal form:
a sorted list of monomials $e_2^p \cdot e_3^q$.
Two expressions are equivalent if and only if
their normal forms coincide.

\section{Cayley--Dickson Compatibility}

Among $\mathbb{R}$, $\mathbb{C}$, $\mathbb{H}$,
$\mathbb{O}$, only $\mathbb{C}$ is compatible
with $\approx$.
$\mathbb{R}$ has a total order absent from $\approx$.
$\mathbb{H}$ lacks commutativity.
$\mathbb{O}$ lacks associativity.
$\mathbb{C}$ satisfies all structural constraints.

\section{Relation to Previous Chapters}

The constants appearing in Chapters~\ref{ch:whyC}--\ref{ch:whyd5}
arise from a single map $\mathrm{eval}: E \to \mathbb{N}$
defined by $e_1 \mapsto 0$, $e_2 \mapsto 2$, $e_3 \mapsto 3$.
Under this map:
\begin{align*}
d &= \mathrm{eval}(e_2) + \mathrm{eval}(e_3) = 5, \\
(n_S, n_T) &= (\mathrm{eval}(e_3), \mathrm{eval}(e_2))
  = (3, 2), \\
K &= \mathbb{C} \quad
  (\text{unique Cayley--Dickson algebra with }
  \dim_{\mathbb{R}} = \mathrm{eval}(e_2) = 2).
\end{align*}
The map $\mathrm{eval}$ is a choice external to the structure.
